\section{Canonical full ambient histories anchored to original normal contact}

Throughout this section $n>0$, $M=(\mathbf C^n,\mathbf C^n,\varnothing,I,1)$
lies in the depth-$q$ invariant domain, and $J=J(I)$ is the full ambient
coherent stabilized record. Let $K$ be its original literal normal contact
record, formed with the original coordinate $t$ on the whole original open
$O=\mathbf C^{n+1}$. Write $R=\operatorname{invFamily}(n+1,q)$ and
$V=\operatorname{invFamily}(n+1,q+1)$. Every tower below starts at the original
full ambient manifold. No current transformed record is restarted.

\begin{lemma}[Persistent upper regimes and full original clocks]
\label{hir:disjoint-canonical-normal-regimes}
\uses{hir:disjoint-full-ambient-unit-monomial-stalk,hir:disjoint-first-step-normal-history,hir:disjoint-coherent-normal-initial-regime}
\lean{TranscendentalNakamaye.Pending.HironakaDisjointCanonicalHistoricalAnchor.canonical_not_special,TranscendentalNakamaye.Pending.HironakaDisjointCanonicalHistoricalAnchor.canonical_order_one,TranscendentalNakamaye.Pending.HironakaDisjointCanonicalHistoricalAnchor.canonical_step_run,TranscendentalNakamaye.Pending.HironakaDisjointCanonicalHistoricalAnchor.canonical_order_run}
\leanok
Let $T$ be any actual finite canonical tower of $V$ for the original $J$.
The original-member list remains empty and its level is zero. The current
codimension-one part is empty, every point is nonspecial, and every point
of the current cosupport has nonmonomial order one. Both original runs at
such a point equal the full actual length of $T$, including empty years.
Their selected heads are the original identity and their tails are $T$.
\end{lemma}
\begin{proof}
The original-member list is the prefix of length $|J.E|=0$ of the current
boundary, so its level is zero regardless of the born members. The natural
initial domain gives maximal order one. Canonicality gives an actual test
tower, so maximal order and absence of codimension-one cosupport persist
along it. A positive boundary exponent at a point on an SNC member would
locally put that smooth hypersurface in the analytic cosupport and produce
a codimension-one point. Hence the monomial-part stalk contains one.
Multiplication by this member in the defining colon makes every germ of
the nonmonomial part vanish at a cosupport point. Its order is positive.
Canonical order monotonicity bounds it by the order one at its original
image, proving exact order one. Empty original members and the empty
codimension-one part exclude both special branches.

For the Step One run use induction on the genuine tower. At every preceding
stage the level-zero nonspecial regime holds, so its recursive run adds one
at each actual blowup, even when its centre is empty. For the order run,
every image of a final cosupport point is a cosupport point at the preceding
canonical prefix. Its order is one by the preceding argument, so this run
also adds one at every actual step. Both initial runs are zero. The actual
split at the full length is the original identity head followed by $T$,
as follows by the recursive split equations, without deleting any year.
\end{proof}

\begin{lemma}[Entire contact cosupport and chronological born boundary]
\label{hir:disjoint-canonical-normal-contact-dictionary}
\uses{hir:disjoint-first-step-literal-contact-canonical,hir:disjoint-canonical-normal-regimes}
\lean{TranscendentalNakamaye.Pending.HironakaDisjointCanonicalHistoricalAnchor.canonical_contact_cosupport,TranscendentalNakamaye.Pending.HironakaDisjointCanonicalHistoricalAnchor.canonical_boundary}
\leanok
The original $K$ lies in the depth-$q$ domain. Along every actual canonical
upper $T$, its transformed cosupport on $T^{-1}O$ is precisely the inverse
image of the transformed cosupport of $J$. Its entire current boundary is
the ordered pullback of the entire current boundary of $J$. The latter has
length $\operatorname{length}(T)$; all empty strict transforms and all born
members retain their original chronological slots.
\end{lemma}
\begin{proof}
The literal coordinate is a genuine tangent direction of the original
coherent mark-one ideal. The natural original domain and empty initial
codimension-one part give its contact domain, with no extra nonzero
assumption. The actual contact-tail theorem, applied to original head the
identity and tail $T$, gives the complete cosupport dictionary on the whole
original open. The original coefficient-homogenized record is equivalent
to $J$, so its transformed cosupport is the same along the true test tower.
At the original stage both whole-open records have the same empty divisor.
Equality of original divisors gives equality of their transformed ordered
divisors along the same restricted tower. Restriction of a genuine marked
test tower pulls back the whole transformed record, so this divisor is
exactly the ordered inverse-image list of the upper divisor. Every blowup
maps every existing slot by its centre-relative strict transform and appends
one inverse-image-centre slot. Starting from the empty list, this gives
length exactly the number of actual steps, without a nonempty-centre count.
\end{proof}

\begin{theorem}[Original contact canonicity and both historical value laws]
\label{hir:disjoint-canonical-normal-history}
\uses{hir:disjoint-canonical-normal-regimes,hir:disjoint-canonical-normal-contact-dictionary,hir:disjoint-first-step-full-historical-anchor}
\lean{TranscendentalNakamaye.Pending.HironakaDisjointCanonicalHistoricalAnchor.canonical_history_data,TranscendentalNakamaye.Pending.HironakaDisjointCanonicalHistoricalAnchor.canonical_history_value}
\leanok
For every actual finite canonical tower $T$ of $V$ for $J$, the same upper
tower is canonical for original Step One, and $T|_O$ is canonical for the
original $K$. At every point $y$ of the transformed upper cosupport,
\[
 \operatorname{StepOne}_R(J,T,y)=(0,R(K,T|_O,y)),\qquad
 V(J,T,y)=\iota_q\bigl(R(K,T|_O,y)\bigr).
\]
Both finite inclusions in $\iota_q$ are retained. The lower canonicity, the
current order, and both clocks are conclusions, not input certificates.
\end{theorem}
\begin{proof}
First fix an upper canonical tower for which both claimed lower canonicalities
have already been proved. The preceding regime lemma identifies the selected
Step One split with original identity head and actual full tail. The genuine
canonical-tail theorem supplies the independently selected lower record's
domain, canonicality and cosupport. The literal original $t$-contact has
its natural domain and its assumed-for-this-intermediate-argument canonical
tail. Since the original members are empty, the level-zero contact anchor
identifies its value with the selected lower value. Nonspecialness gives
the first displayed formula.

For the outer invariant its original order run also selects identity head
and actual full tail, with order one. The original order-one mask is the
whole original open. The original empty-boundary companion is restricted-
equivalent to the whole-open restriction of $J$. Its domain and maximal
order are natural companion consequences. Restriction of the original
canonical Step One tail and representative invariance make its actual
companion tail canonical and give the same Step One value. The outer
positive-order constructor therefore gives exactly the second formula.
These intermediate arguments do not construct either missing lower
canonicality from its own future value.

Now prove both canonicalities and both value formulas simultaneously by
induction on $T$. For the identity tower all canonicalities are immediate,
so the two intermediate arguments prove the two initial formulas. Suppose
all four conclusions hold for the genuine canonical prefix $P$, and append
an actual blowup with centre $D$. If $D$ is empty, its restricted centre is
empty, so both canonicalities extend. Otherwise $D$ is the entire maximum
of the upper value at $P$. On its whole cosupport, the induction hypothesis
writes the upper values as $\iota_q$ of the actual original contact values,
and Step One values as $(0,\text{finite contact value})$. Both constructors
are order embeddings. Thus the maximum of upper values is exactly the
maximum of Step One values. The whole contact cosupport dictionary also
identifies its restricted centre with the entire maximum of the original
contact value along $P|_O$. This proves both next canonicalities using only
the already proved prefix formulas. Apply the intermediate Step One and
outer arguments to the extended tower, obtaining both next formulas.
The induction is on actual tower length and does not restart at the endpoint
or erase empty years. Its conclusions give the theorem for every finite $T$.
\end{proof}

\begin{theorem}[Every full ambient next maximum is the original contact maximum]
\label{hir:disjoint-canonical-normal-maximum}
\uses{hir:disjoint-canonical-normal-history,hir:disjoint-canonical-normal-contact-dictionary}
\lean{TranscendentalNakamaye.Pending.HironakaDisjointCanonicalHistoricalAnchor.canonical_history_maximum_locus}
\leanok
At every original upper canonical prefix, the full ambient next maximum
locus equals the image under the actual whole-open inclusion of the maximum
locus of the original literal contact invariant along its actual restricted
history. This includes empty transformed cosupports and arbitrary empty years.
\end{theorem}
\begin{proof}
The whole transformed cosupport dictionary gives the same points on both
sides, and every upper point has its unique whole-open lift. The preceding
historical value formula identifies each full ambient value with $\iota_q$
of that lift's contact value. Since $\iota_q$ preserves and reflects order,
an upper point dominates all upper cosupport points if and only if its
whole-open lift dominates all contact cosupport points. This proves both
set inclusions. If either transformed cosupport is empty, the same
dictionary makes both empty, and both maximum loci are empty by definition.
\end{proof}

\begin{remark}[Remaining cross-ambient comparison]
The contact on the right starts from the original thin carrier $H$ and keeps
its entire actual history, born boundary and original references. It is not
the invariant of a fresh current record, and is not identified here with the
unmodified raw base invariant. A further thin-to-base historical comparison
must account for the same-index ambient-dimension map in nonbottom monomial
branches. No such comparison or full intrinsic stabilization descent is
part of these theorems.
\end{remark}
